% ex: ts=2 sw=2 sts=2 et filetype=tex
% SPDX-License-Identifier: CC-BY-SA-4.0

\section{Introducción}

\begin{frame}[c]{Introducción}

En la actualidad, la educación enfrenta el reto de adaptarse a
entornos digitales cada vez más dinámicos, donde el uso de
\textbf{Tecnologías de la Información y la Comunicación} (TIC) es
fundamental para mejorar los procesos de \emph{enseñanza-aprendizaje}.

\vspace{1em}
La evolución tecnológica ha desplazado los métodos tradicionales
(como la oralidad y los textos impresos) hacia la virtualidad.
\end{frame}

\section{Trabajos relacionados}

\begin{frame}[c]{Trabajos relacionados}

Plataformas como \textbf{Google Classroom}, \textbf{Moodle} y
los \emph{cursos en línea masivos y abiertos} o \textbf{MOOC} (acrónimo
en inglés de \emph{massive open online courses}) han permitido
ampliar el acceso a recursos educativos; sin embargo, aún existen
desafíos importantes en la personalización del aprendizaje y la
evaluación continua.

  \vspace{1em}
  \begin{center}
    
\includegraphics[height=1.5cm]{g-classroom.png}
    \hspace{1em}
    
\includegraphics[height=1.5cm]{g-forms.png}
    \hspace{1em}
    
\includegraphics[height=1.5cm]{ms-forms.png}
    \hspace{1em}
    
\includegraphics[height=1.5cm]{canvas.png}
    \hspace{1em}
    
\includegraphics[height=1.5cm]{moodle.png}
  \end{center}
\end{frame}

\begin{frame}[c]{Áreas de mejora}
  \begin{itemize}
    \item Los sistemas actuales \textbf{no} generan preguntas
        a\textbf{utomáticas} o \textbf{aleatorias}.
    \pausa
    \item Las preguntas \textbf{casi siempre} son las mismas para
        todos los estudiantes.
    \begin{itemize}
      \item Se pueden copiar las respuesta entre los alumnos
      \item Se pueden copiar las preguntas y guardarlas para
          el siguiente ciclo.
    \end{itemize}
    \pausa
    \item Se necesita tener \textbf{acceso a Internet} para todos
        los alumnos.
    \begin{itemize}
      \item Uso de PC o en rara vez en ceular.
    \end{itemize}
  \end{itemize}
\end{frame}

\section{Justificación}

\begin{frame}[c]{Justificación}
  \begin{itemize}
    \item Crear un sistema que genere preguntas \textbf{automáticas} o
        \textbf{aleatorias} de matemáticas.
    \pausa
    \item \textbf{Evitar la copia} entre alumnos al ser preguntas
        diferentes.
    \pausa
    \item Acceso al sistema \textbf{sin necesidad} de Internet.
  \end{itemize}
\end{frame}

\section{Objetivos}

\begin{frame}[c]{Objetivo general}
  \begin{itemize}
    \item Desarrollar un asistente virtual que permita la generación
      y evaluación automática de ejercicios matemáticos para apoyar
      la adquisición de competencias técnicas en estudiantes de
      educación media superior.
  \end{itemize}
\end{frame}

\begin{frame}[c]{Objetivos específicos}
  \begin{enumerate}
    \item \label{itm:objetivo1} Desarrollar un sistema que realice exámenes
      y ejercicios de manera autónoma y aleatoria.
    \pausa
    \item \label{itm:objetivo2} Crear una herramienta que permita a los
      profesores crear, modificar y administrar las pruebas y ejercicios de
      aprendizaje.
    \pausa
    \item \label{itm:objetivo3} Crear una herramienta que permita a los
      alumnos realizar sus exámenes y ejercicios.
    \pausa
    \item \label{itm:objetivo4} Crear un ambiente seguro para la realización
      de los exámenes.
    \pausa
    \item \label{itm:objetivo5} Crear un sistema que no requiera la conexión
      a Internet para que se conecten los alumnos.
  \end{enumerate}
\end{frame}

\section{Marco Teórico}

\begin{frame}[c]{Marco Teórico}
  \begin{enumerate}
    \item Uso del framework Laravel
      \begin{enumerate}
        \item Usando el lenguaje PHP
      \end{enumerate}
    \pausa
    \item Uso de JavaScript para hacer ejercicios del lado del cliente
      \begin{enumerate}
        \item \textbf{MathJax} para el despliegue de LaTeX, MathML y
          AsciiMath
      \end{enumerate}
    \pausa
    \item Uso de una Raspberry Pi o una NUC como servidor de red local
  \end{enumerate}

  \vspace{1em}
  \begin{center}
    \includegraphics[height=2cm]{rpi4.png}
    \hspace{3em}
    \includegraphics[height=2cm]{nuc.png}
  \end{center}
\end{frame}

\section{Propuesta}

\begin{frame}[c]{Propuesta}
  \begin{enumerate}
    \item Creación de una aplicación web que no necesite conexión a
      Internet.
    \pausa
    \item Uso de Hardware de bajo consumo como la Raspberry Pi o la NUC para
      actuar como servidor local.
      \begin{itemize}
        \item El tamaño pequeño ayuda a su transportación y manejo en el
          salón.
        \item Usar este mismo equipo como punto de acceso inalámbrico
          (un hostpot), para acceder al sistema.
      \end{itemize}
  \end{enumerate}
\end{frame}

\section{Resultados}

\begin{frame}[c]{Pantalla de inicio}
  \begin{columns}
    \column{0.3\textwidth}
      \begin{enumerate}
        \item Se creo un sistema responsivo (para PC y moviles)
        \item Con la característica de utilizar dos estilos: claro y oscuro
      \end{enumerate}
    \column{0.7\textwidth}
      \begin{center}
        \includegraphics[width=5cm]{app/01-inicio.png}
        \includegraphics[width=5cm]{app/02-inicio.png}
      \end{center}
  \end{columns}
\end{frame}

\begin{frame}[c]{Pantalla de registro}
  \begin{columns}
    \column{0.3\textwidth}
      \begin{enumerate}
        \item Capacidad de registrar nuevos usuarios
      \end{enumerate}
    \column{0.7\textwidth}
      \begin{center}
        \includegraphics[width=4.5cm]{app/03-registro.png}
        \includegraphics[width=4.5cm]{app/04-registro.png}
      \end{center}
  \end{columns}
\end{frame}

\begin{frame}[c]{Pantalla de inicio de sesión}
  \begin{columns}
    \column{0.5\textwidth}
      \begin{enumerate}
        \item Capacidad de iniciar sesión con usuario o email
      \end{enumerate}
    \column{0.5\textwidth}
      \begin{center}
        \includegraphics[width=6cm]{app/05-isesion.png}
      \end{center}
  \end{columns}
\end{frame}

\begin{frame}[c]{Pantalla de Panel principal alumno}
  \begin{columns}
    \column{0.5\textwidth}
      \begin{enumerate}
        \item Al iniciar el alumno vera sus materias/cursos
            al que esta inscrito
      \end{enumerate}
    \column{0.5\textwidth}
      \begin{center}
        \includegraphics[width=6cm]{app/07-dashboard.png}
      \end{center}
  \end{columns}
\end{frame}

\begin{frame}[c]{Pantalla de Panel principal administrador}
  \begin{columns}
    \column{0.5\textwidth}
      \begin{enumerate}
        \item El administrador podrá manipular los usuarios y
          las asignaturas
      \end{enumerate}
    \column{0.5\textwidth}
      \begin{center}
        \includegraphics[width=6cm]{app/06-dashboard.png}
      \end{center}
  \end{columns}
\end{frame}

\begin{frame}[c]{Pantalla de usuarios}
  \begin{columns}
    \column{0.5\textwidth}
      \begin{enumerate}
        \item Lista de usuarios y la capacidad de modificarlos
      \end{enumerate}
    \column{0.5\textwidth}
      \begin{center}
        \includegraphics[width=6cm]{app/08-usuarios.png}
      \end{center}
  \end{columns}
\end{frame}

\begin{frame}[c]{Pantalla de asignaturas}
  \begin{columns}
    \column{0.5\textwidth}
      \begin{enumerate}
        \item Lista de asignaturas y su manipulación
      \end{enumerate}
    \column{0.5\textwidth}
      \begin{center}
        \includegraphics[width=6cm]{app/09-asignaturas.png}
      \end{center}
  \end{columns}
\end{frame}

\begin{frame}[c]{Pantalla de grupos de las asignaturas}
  \begin{columns}
    \column{0.5\textwidth}
      \begin{enumerate}
        \item Cada asignatura tiene sus grupos
        \item Cada grupo tiene un maestro asignado
      \end{enumerate}
    \column{0.5\textwidth}
      \begin{center}
        \includegraphics[width=7cm]{app/10-algebra.png}
      \end{center}
  \end{columns}
\end{frame}

\begin{frame}[c]{Pantalla de alumnos registrados}
  \begin{columns}
    \column{0.5\textwidth}
      \begin{enumerate}
        \item Cada grupo tiene sus alumnos
      \end{enumerate}
    \column{0.5\textwidth}
      \begin{center}
        \includegraphics[width=6cm]{app/11-grupo17.png}
      \end{center}
  \end{columns}
\end{frame}

\begin{frame}[c]{Pantalla de ejemplo de ejercicio}
  \begin{columns}
    \column{0.5\textwidth}
      \begin{enumerate}
        \item Ejercicio de resolver una ecuación cuadrática
      \end{enumerate}
    \column{0.5\textwidth}
      \begin{center}
        \includegraphics[width=4.5cm]{app/12-ecua-cuadra.png}
      \end{center}
  \end{columns}
\end{frame}

\begin{frame}[c]{Pantalla de ejemplo de ejercicio}
  \begin{columns}
    \column{0.5\textwidth}
      \begin{enumerate}
        \item Mensaje de error
      \end{enumerate}
    \column{0.5\textwidth}
      \begin{center}
        \includegraphics[width=4.5cm]{app/13-ecua-cuadra.png}
      \end{center}
  \end{columns}
\end{frame}

\begin{frame}[c]{Pantalla de ejemplo de ejercicio}
  \begin{columns}
    \column{0.5\textwidth}
      \begin{enumerate}
        \item Mensaje de éxito
      \end{enumerate}
    \column{0.5\textwidth}
      \begin{center}
        \includegraphics[width=4.5cm]{app/14-ecua-cuadra.png}
      \end{center}
  \end{columns}
\end{frame}

\section{Conclusiones}

\begin{frame}[c]{Conclusiones}
Este sistema nos da las siguientes ventajas:
  \begin{enumerate}
    \item Los alumnos \textbf{no están conectados a Internet}, para usar la
      IA o google para hacer trampa.
  \pausa
  \item \textbf{No se necesita Internet} para usar el sistema
  \pausa
  \item Los \textbf{ejercicios} y \textbf{exámenes} se generan y califican
    \textbf{automáticamente}
  \pausa
  \item Las preguntas son \textbf{diferentes} para cada alumno
  \pausa
  \item Los alumnos pueden usar una PC o su propio celular
  \end{enumerate}
\end{frame}

\section{Trabajo futuro}

\begin{frame}[c]{Trabajo futuro}
  \begin{enumerate}
    \item Exportar las calificaciones a Excel o CSV.
    \item Hacer respaldos en la nube cuando se requiera.
    \item Conectarlos con alguna otra plataforma como Google Classroom.
    \item Importar nuevas preguntas.
    \item Ademas de JavaScrilpt utilizar otros backends.
    \item Alojarlo en un sitio web en internet para que haga ejercicios.
  \end{enumerate}
\end{frame}

\begin{frame}[c]{Gracias}
  ¿Preguntas?
  \begin{itemize}
    \item El código se encuentra en: https://github.com/yolandayah/emilie
  \end{itemize}
\end{frame}
